%% =============================================================================
%% chapters/appendix-praxis-workflow.tex
%% Appendix D: Implementation and Deployment Guide for Praxis Workflow Sequences
%% =============================================================================

\chapter{Implementation and Deployment Guide for Praxis Workflow Sequences}
\label{app:praxis_workflow}

\section{Overview \& Engineering Objectives}
\label{sec:praxis_guide_overview}

Within the Harmonia Health Information Exchange (HIE) architecture, the \textbf{Energeia} subsystem governs all workflow orchestration, clinical transformation, data enrichment, and task execution. While individual tasks are encapsulated within discrete, modular \textbf{Ergon} activity modules (documented in Appendix~\ref{app:ergon_modules}), the deterministic composition, ordering, topic-based dispatch, dynamic Apache Camel route compilation, and end-to-end audit checkpointing of these activities are executed by \textbf{Praxis Workflow Sequences}.

This implementation guide provides an authoritative, production-grade technical manual for software engineers, integration developers, and platform architects charged with authoring, configuring, testing, and deploying Praxis workflow sequences. Specifically, this guide covers:
\begin{enumerate}
    \item \textbf{The Praxis Architectural Paradigm}: Decoupling declarative workflow metadata (\texttt{PraxisDefinition} in the \texttt{calliope} model module) from runtime Camel execution (\texttt{PraxisImplementation} and \texttt{Praxis} in the \texttt{energeia/praxis} engine), managing multi-level hierarchical topic subscriptions (\texttt{TopicSubscription}), configuring target gateway routing, and persisting execution state across a dual-cache Infinispan grid.
    \item \textbf{Step-by-Step Praxis Creation Blueprint}: Defining declarative sequence descriptors in JSON and Java, ordering pipeline activities, binding ingress queues and endpoints, and managing error containment boundaries.
    \item \textbf{Production Reference Implementation}: A complete, production-ready implementation of \texttt{PatientCareAdmissionPraxis}, alongside its companion declarative JSON configuration, demonstrating multi-step clinical orchestration spanning patient demographics ingestion, identifier cross-referencing, and clinical encounter enrichment.
    \item \textbf{Dynamic Route Generation \& Camel Pipeline Topology}: In-depth analysis of \texttt{createSequencePipelineRoute()}, dynamic direct endpoint synthesis (\texttt{direct:seq-<praxisId>-start} $\rightarrow$ \texttt{direct:seq-<praxisId>-step-N} $\rightarrow$ \texttt{direct:seq-<praxisId>-complete}), ingress conduit bridging, and exchange property invariants.
    \item \textbf{Dynamic Discovery, Seeding \& CDI Runtime Binding}: Operating \texttt{TaskSequenceLoader} to dynamically query \texttt{PraxisService}, discovering \texttt{@Dependent} \texttt{ErgonBase} CDI beans via Jakarta EE \texttt{Instance<ErgonBase>}, bootstrapping default workflows via \texttt{TaskSequenceDefaultSeeder}, and injecting active routes into the running \texttt{CamelContext}.
    \item \textbf{Multi-Tier Testing Strategies}: Authoring deterministic unit and integration test harnesses using Apache Camel Test Kit (\texttt{CamelTestSupport}), synthetic \texttt{Pragma} envelopes, mock endpoint assertions, and verification of transactional \texttt{PragmaCheckpoint} audit trails.
    \item \textbf{Deployment, Cluster Management \& Operational Runbook}: Containerizing within the Ponos WildFly execution environment, monitoring cluster state via the \texttt{/workflow} REST administrative API, executing zero-downtime dynamic sequence reloads via \texttt{ponos-cli --reload-sequences}, and managing Infinispan cache replication.
\end{enumerate}

% =============================================================================
% SECTION D.1: ARCHITECTURAL BLUEPRINT & THE PRAXIS PARADIGM
% =============================================================================
\section{Architectural Blueprint \& the Praxis Paradigm}
\label{sec:praxis_paradigm}

\begin{figure}[htbp]
    \centering
    \begin{adjustbox}{max width=\textwidth}
        %% =============================================================================
%% diagrams/fig-praxis-workflow-architecture.tex
%% ArchiMate 3.2 Praxis Workflow Sequence Architecture & Pipeline Synthesis View
%% =============================================================================
\begin{tikzpicture}[node distance=1.0cm and 0.8cm, auto]

    % --- Column 1: Declarative Metadata & Seeding Tier ---
    \node[archi-data-object] (praxis_def_json) at (0, 7.5) {%
        \archibadge{Data Object}\archiname{Praxis Definition}\\\archidesc{JSON Descriptor / Config}%
    };

    \node[archi-app-component] (praxis_service) at (0, 4.5) {%
        \archibadge{App Component}\archiname{Praxis Service}\\\archidesc{PraxisService HotRod Client}%
    };

    \node[archi-node] (sequence_cache) at (0, 1.5) {%
        \archibadge{Technology Node}\archiname{Task Sequence Cache}\\\archidesc{tasksequence-cache Grid}%
    };

    % --- Column 2: Dynamic Discovery & CDI Injection Tier ---
    \node[archi-app-component] (sequence_loader) at (4.5, 7.5) {%
        \archibadge{App Component}\archiname{Sequence Loader}\\\archidesc{TaskSequenceLoader CDI}%
    };

    \node[archi-app-component] (cdi_ergon_beans) at (4.5, 4.5) {%
        \archibadge{App Component}\archiname{CDI Ergon Resolution}\\\archidesc{Instance<ErgonBase> Injection}%
    };

    \node[archi-app-component] (praxis_impl) at (4.5, 1.5) {%
        \archibadge{App Component}\archiname{Praxis Implementation}\\\archidesc{Praxis / Route Builder}%
    };

    % --- Column 3: Dynamic Camel Pipeline Topology Tier ---
    \node[archi-app-interface] (pipeline_start) at (9.0, 7.5) {%
        \archibadge{App Interface}\archiname{Pipeline Ingress}\\\archidesc{direct:seq-<praxisId>-start}%
    };

    \node[archi-app-process] (pipeline_steps) at (9.0, 4.5) {%
        \archibadge{App Process}\archiname{Chained Erga Steps}\\\archidesc{direct:seq-step-0 \dots\ step-N}%
    };

    \node[archi-app-interface] (pipeline_complete) at (9.0, 1.5) {%
        \archibadge{App Interface}\archiname{Pipeline Egress}\\\archidesc{direct:seq-<praxisId>-complete}%
    };

    % --- Column 4: Checkpointing, Storage & Runtime Operations Tier ---
    \node[archi-app-component] (checkpoint_manager) at (13.5, 7.5) {%
        \archibadge{App Component}\archiname{Checkpoint Manager}\\\archidesc{PraxisCheckpointManager}%
    };

    \node[archi-node] (pragma_cache) at (13.5, 4.5) {%
        \archibadge{Technology Node}\archiname{Pragma Audit Cache}\\\archidesc{pragma-cache Infinispan}%
    };

    \node[archi-app-component] (mgmt_resource) at (13.5, 1.5) {%
        \archibadge{App Component}\archiname{Workflow Ops REST}\\\archidesc{/workflow \& ponos-cli}%
    };

    % --- Flow & Structural Relationships ---
    % Declarative loading and cache
    \draw[archi-access] (praxis_service) -- (praxis_def_json);
    \draw[archi-flow] (praxis_service) -- node[right, font=\tiny] {HotRod CRUD} (sequence_cache);
    \draw[archi-serving] (praxis_service) -- (sequence_loader);

    % Dynamic resolution and instantiation
    \draw[archi-triggering] (sequence_loader) -- (cdi_ergon_beans);
    \draw[archi-assignment] (cdi_ergon_beans) -- (praxis_impl);
    \draw[archi-triggering] (sequence_loader) -- (praxis_impl);

    % Dynamic Camel Route Generation
    \draw[archi-triggering] (praxis_impl) -- node[above, font=\tiny] {createSequenceRoute} (pipeline_start);
    \draw[archi-flow] (pipeline_start) -- (pipeline_steps);
    \draw[archi-flow] (pipeline_steps) -- (pipeline_complete);

    % Checkpointing and State Persistence
    \draw[archi-flow] (pipeline_steps) -- node[above, font=\tiny] {Audit Events} (checkpoint_manager);
    \draw[archi-flow] (checkpoint_manager) -- node[right, font=\tiny] {Save Checkpoint} (pragma_cache);

    % Operations & Dynamic Reload
    \draw[archi-serving] (mgmt_resource) -- (sequence_loader);
    \draw[archi-access] (mgmt_resource) -- (sequence_cache);

\end{tikzpicture}

    \end{adjustbox}
    \caption{ArchiMate 3.2 Praxis Workflow Sequence Architecture and Dynamic Route Synthesis Topology}
    \label{fig:praxis_workflow_architecture}
\end{figure}

The Harmonia workflow engine enforces a strict architectural separation of concerns between declarative workflow definition and runtime route execution. Figure~\ref{fig:praxis_workflow_architecture} illustrates the ArchiMate 3.2 processing and component topology of the Praxis workflow subsystem.

\subsection{Decoupling Definition from Runtime Execution}
\label{subsec:praxis_decoupling}

Workflow sequences are architected into two distinct layers:
\begin{itemize}
    \item \textbf{Declarative Metadata Layer (\texttt{PraxisDefinition})}: Residing in the shared \texttt{calliope} model module (\texttt{net.fhirfactory.harmonia.model.praxis}), \texttt{PraxisDefinition} is a lightweight, serializable Java Plain Old Java Object (POJO) and JSON data transfer object. It encapsulates workflow metadata, including the unique sequence identifier (\texttt{praxisId}), human-readable name, semantic version, enabled flag, topic subscriptions (\texttt{List<TopicSubscription>}), target gateway instance IDs (\texttt{List<String>}), trigger types (\texttt{List<String>}), and an ordered map of activity class names (\texttt{Map<Integer, String> activityClassNames}). It possesses no runtime dependencies on Apache Camel, CDI, or container runtimes, enabling safe network transport and JSON persistence in the Infinispan cache.
    \item \textbf{Runtime Execution Layer (\texttt{PraxisImplementation} / \texttt{Praxis})}: Residing in the \texttt{energeia/praxis} engine (\texttt{net.fhirfactory.harmonia.praxis.sequence}), \texttt{PraxisImplementation} extends \texttt{PraxisDefinition} and acts as an Apache Camel \texttt{RouteBuilder}. It dynamically resolves concrete \texttt{@Dependent} \texttt{ErgonBase} activity beans from the CDI container, binds transactional \texttt{PraxisCheckpointManager} interceptors, and compiles the dynamic pipeline routing logic into the active \texttt{CamelContext}.
\end{itemize}

\subsection{Topic Subscription Matching \& Hierarchical Routing}
\label{subsec:praxis_topic_subscriptions}

Praxis workflows register topic subscriptions that govern which incoming clinical events trigger pipeline execution. A \texttt{Topic} in Harmonia consists of a structured four-tuple:
\begin{equation}
\text{Topic} = \langle \text{Domain}, \text{Model}, \text{Element}, \text{Qualifier} \rangle
\end{equation}

For example, an inpatient admission ADT A01 message from a Patient Administration System maps to the topic:
\begin{lstlisting}[caption={Hierarchical Topic Representation}]
Domain:    HIE
Model:     HL7_V2
Element:   ADT_A01
Qualifier: INPATIENT_ADMISSION
\end{lstlisting}

A \texttt{TopicSubscription} defines matching rules with wildcard capabilities (\texttt{*}):
\begin{itemize}
    \item \textbf{Exact Match}: Matches a specific topic tuple exactly (e.g., \texttt{HIE/HL7\_V2/ADT\_A01/INPATIENT\_ADMISSION}).
    \item \textbf{Wildcard Subscriptions}: Matches all elements under a model or domain (e.g., \texttt{HIE/HL7\_V2/ADT\_*/\allowbreak *}).
    \item \textbf{Priority Resolution}: When an incoming \texttt{Pragma} matches multiple registered Praxis sequences, the dispatcher routes the task based on explicit sequence priority weights or target gateway instance affinity.
\end{itemize}

\subsection{Audit Checkpoint Lifecycle across Execution Boundaries}
\label{subsec:praxis_checkpoint_lifecycle}

To guarantee non-repudiation, HIPAA audit compliance, and automated failure recovery, the \texttt{PraxisCheckpointManager} records durable audit checkpoints at every major execution transition. Table~\ref{tab:praxis_checkpoint_stages} details the five canonical checkpoint stages.

\begin{table}[htbp]
\centering
\small
\begin{tabularx}{\textwidth}{l p{3.2cm} X l}
\toprule
\textbf{Stage Constant} & \textbf{Stage Name} & \textbf{Description \& Audit State Recorded} & \textbf{Pragma Status} \\
\midrule
\texttt{STAGE\_PIPELINE\_INGRESS} & \texttt{PIPELINE\_INGRESS} & Ingress handoff received at \texttt{direct:seq-<id>-start}; records initial payload hash and ingress timestamp. & \texttt{QUEUED} \\
\texttt{STAGE\_PRE\_ERGON} & \texttt{PRE\_ERGON} & Intercepted immediately before delegating to \texttt{ErgonBase} step $N$; records step index and input state snapshot. & \texttt{PROCESSING} \\
\texttt{STAGE\_POST\_ERGON} & \texttt{POST\_ERGON} & Intercepted upon successful return from \texttt{ErgonBase} step $N$; records transformed FHIR resources and execution duration. & \texttt{PROCESSING} \\
\texttt{STAGE\_PIPELINE\_COMPLETE} & \texttt{PIPELINE\_COMPLETE} & Pipeline egress reached at \texttt{direct:seq-<id>-complete}; records total execution time and final outcome envelope. & \texttt{COMPLETED} \\
\texttt{STAGE\_PIPELINE\_FAILED} & \texttt{PIPELINE\_FAILED} & Trapped by pipeline error boundary; records failure reason, exception class, stack trace, and failed step index. & \texttt{FAILED} \\
\bottomrule
\end{tabularx}
\caption{Praxis Workflow Execution Checkpoint Stages and Lifecycle Semantics}
\label{tab:praxis_checkpoint_stages}
\end{table}

\subsection{Dual-Cache Architecture in the Mneme Grid}
\label{subsec:praxis_dual_cache}

State persistence in the \texttt{Mneme} distributed Infinispan grid is partitioned into two specialized caches:
\begin{enumerate}
    \item \textbf{\texttt{tasksequence-cache}}: Stores declarative \texttt{PraxisDefinition} JSON configurations. Accessed via \texttt{PraxisService}, this cache is replicated across all Ponos cluster nodes. Any update to this cache triggers a cluster-wide cache invalidation and hot-reloads the active Camel routing topology.
    \item \textbf{\texttt{pragma-cache}}: Stores active and completed in-flight \texttt{Pragma} task envelopes and their associated \texttt{PragmaCheckpoint} historical audit chains. Checkpoints written by \texttt{PraxisCheckpointManager} are synchronously committed using the HotRod binary protocol to guarantee durability before subsequent pipeline steps execute.
\end{enumerate}

% =============================================================================
% SECTION D.2: STEP-BY-STEP PRAXIS CREATION BLUEPRINT
% =============================================================================
\section{Step-by-Step Praxis Creation Blueprint}
\label{sec:praxis_creation_blueprint}

Engineering a new Praxis workflow sequence follows a structured four-phase lifecycle: definition, activity composition, route configuration, and error boundary definition.

\subsection{Phase 1: Defining Declarative Workflow Metadata}
\label{subsec:blueprint_metadata}

A Praxis workflow sequence must define unique identification, human-readable descriptors, topic subscriptions, and target gateways. In declarative environments (e.g., Kubernetes ConfigMaps or Infinispan seeders), this metadata is expressed as JSON:

\begin{lstlisting}[caption={Declarative JSON Descriptor for Praxis Workflow Sequence}]
{
  "praxisId": "praxis-patient-admission",
  "name": "Patient Care Admission Workflow",
  "description": "Orchestrates patient admission ingestion, demographic parsing, identifier lookup, and clinical encounter generation.",
  "version": "1.0.0",
  "enabled": true,
  "topicSubscriptions": [
    {
      "topic": {
        "domain": "HIE",
        "model": "HL7_V2",
        "element": "ADT_A01",
        "qualifier": "INPATIENT"
      },
      "active": true
    }
  ],
  "targetGateways": ["gateway-pylai-mllp-in-pas"],
  "targetTriggers": ["INBOUND_HL7_EVENT"],
  "sourceQueueName": "petasos.queue.inbound.adt",
  "activityClassNames": {
    "0": "net.fhirfactory.harmonia.erga.patient.demographics.PatientDemographicsUpdateErgon",
    "1": "net.fhirfactory.harmonia.erga.patient.identity.PatientIdentifierCrossReferenceErgon",
    "2": "net.fhirfactory.harmonia.erga.encounter.EncounterEnrichmentErgon"
  }
}
\end{lstlisting}

\subsection{Phase 2: Activity Ordering and Sequence Validation}
\label{subsec:blueprint_validation}

When instantiating a \texttt{PraxisImplementation}, the runtime executes the \texttt{validate()} contract to enforce topological invariants:
\begin{enumerate}
    \item \textbf{Sequence Identifier Validation}: Ensures \texttt{praxisId} is non-null and matches the naming pattern \verb|/^[a-z0-9-]+$/|.
    \item \textbf{Zero-Gapped Activity Indexing}: Ensures activities are indexed consecutively starting from index \texttt{0} through $N-1$.
    \item \textbf{Bean Resolution Non-Null Guarantee}: Validates that every configured activity class name successfully resolves to an active, non-null \texttt{ErgonBase} instance in the CDI container.
\end{enumerate}

\subsection{Phase 3: Route Synthesis and Ingress Conduits}
\label{subsec:blueprint_ingress}

The \texttt{PraxisImplementation} automatically generates two primary Camel route structures:
\begin{itemize}
    \item \textbf{Ingress Bridge Route}: If \texttt{sourceQueueName} is configured, an ingress JMS route consumes from \texttt{jms:queue:<sourceQueueName>} and forwards exchanges to \texttt{direct:seq-<praxisId>-start}.
    \item \textbf{Linear Activity Pipeline}: Synthesizes a chained Camel route connecting \texttt{direct:seq-<praxisId>-start} through intermediate activity endpoints \texttt{direct:seq-<praxisId>-step-N} to the completion sink \texttt{direct:seq-<praxisId>-complete}.
\end{itemize}

\subsection{Phase 4: Error Boundary \& Exception Trapping}
\label{subsec:blueprint_error_handling}

The pipeline route establishes an explicit \texttt{onException(Exception.class)} error trap:
\begin{lstlisting}[language=java, caption={Camel Error Boundary in Praxis Route Synthesis}]
onException(Exception.class)
    .handled(true)
    .process(exchange -> {
        Exception cause = exchange.getProperty(Exchange.EXCEPTION_CAUGHT, Exception.class);
        Pragma pragma = resolvePragma(exchange);
        if (pragma != null && checkpointManager != null) {
            checkpointManager.recordPipelineFailure(pragma, getPraxisId(), currentStepIndex, cause);
        }
        exchange.getMessage().setHeader("PraxisExecutionStatus", "FAILED");
    })
    .to("direct:seq-" + getPraxisId() + "-error")
    .to("jms:queue:petasos.queue.dlq.workflow-failure");
\end{lstlisting}

% =============================================================================
% SECTION D.3: PRODUCTION REFERENCE IMPLEMENTATION
% =============================================================================
\section{Production Reference Implementation: PatientCareAdmissionPraxis}
\label{sec:praxis_reference_impl}

Listing~\ref{lst:patient_care_admission_praxis} provides the complete, production-grade Java reference implementation for \texttt{PatientCareAdmissionPraxis}.

\begin{lstlisting}[language=java, caption={Production Reference Implementation: PatientCareAdmissionPraxis.java}, label={lst:patient_care_admission_praxis}]
package net.fhirfactory.harmonia.praxis.clinical;

import jakarta.enterprise.context.Dependent;
import jakarta.inject.Inject;
import net.fhirfactory.harmonia.erga.base.ErgonBase;
import net.fhirfactory.harmonia.erga.encounter.EncounterEnrichmentErgon;
import net.fhirfactory.harmonia.erga.patient.demographics.PatientDemographicsUpdateErgon;
import net.fhirfactory.harmonia.erga.patient.identity.PatientIdentifierCrossReferenceErgon;
import net.fhirfactory.harmonia.model.praxis.Topic;
import net.fhirfactory.harmonia.model.praxis.TopicSubscription;
import net.fhirfactory.harmonia.praxis.sequence.Praxis;
import org.apache.camel.CamelContext;
import org.slf4j.Logger;
import org.slf4j.LoggerFactory;

import java.util.ArrayList;
import java.util.HashMap;
import java.util.List;
import java.util.Map;

/**
 * Production reference implementation of a multi-step clinical admission workflow.
 * Orchestrates:
 *   Step 0: PatientDemographicsUpdateErgon (Extracts patient demographic FHIR Patient)
 *   Step 1: PatientIdentifierCrossReferenceErgon (Queries EMPI / Master Patient Index)
 *   Step 2: EncounterEnrichmentErgon (Constructs FHIR Encounter and location mappings)
 */
@Dependent
public class PatientCareAdmissionPraxis extends Praxis {

    private static final Logger LOG = LoggerFactory.getLogger(PatientCareAdmissionPraxis.class);

    public static final String PRAXIS_ID = "praxis-patient-care-admission";
    public static final String PRAXIS_NAME = "Patient Care Admission Workflow";
    public static final String PRAXIS_VERSION = "1.2.0";

    @Inject
    public PatientCareAdmissionPraxis(
            CamelContext camelContext,
            PatientDemographicsUpdateErgon demographicsErgon,
            PatientIdentifierCrossReferenceErgon identityErgon,
            EncounterEnrichmentErgon encounterErgon) {
        
        super();
        setCamelContext(camelContext);
        setPraxisId(PRAXIS_ID);
        setName(PRAXIS_NAME);
        setVersion(PRAXIS_VERSION);
        setEnabled(true);
        setDescription("Production multi-step admission orchestration pipeline.");

        // 1. Configure Ingress Queue Binding
        setSourceQueueName("petasos.queue.inbound.adt");
        setInputEndpoint("direct:seq-" + PRAXIS_ID + "-start");

        // 2. Configure Topic Subscriptions
        List<TopicSubscription> subscriptions = new ArrayList<>();
        Topic adtA01Topic = new Topic("HIE", "HL7_V2", "ADT_A01", "INPATIENT");
        subscriptions.add(new TopicSubscription(adtA01Topic, true));
        Topic adtA04Topic = new Topic("HIE", "HL7_V2", "ADT_A04", "EMERGENCY");
        subscriptions.add(new TopicSubscription(adtA04Topic, true));
        setTopicSubscriptions(subscriptions);

        // 3. Configure Target Gateways and Triggers
        List<String> gateways = new ArrayList<>();
        gateways.add("gateway-pylai-mllp-in-pas");
        setTargetGateways(gateways);

        List<String> triggers = new ArrayList<>();
        triggers.add("INBOUND_HL7_EVENT");
        setTargetTriggers(triggers);

        // 4. Bind Sequential Ergon Activities
        Map<Integer, ErgonBase> activityMap = new HashMap<>();
        activityMap.put(0, demographicsErgon);
        activityMap.put(1, identityErgon);
        activityMap.put(2, encounterErgon);
        setActivities(activityMap);

        // 5. Populate Declarative Metadata Map for Serialization
        Map<Integer, String> classNames = new HashMap<>();
        classNames.put(0, demographicsErgon.getClass().getName());
        classNames.put(1, identityErgon.getClass().getName());
        classNames.put(2, encounterErgon.getClass().getName());
        setActivityClassNames(classNames);

        LOG.info("PatientCareAdmissionPraxis initialized with 3 sequential Erga activities.");
    }

    /**
     * Default constructor for dynamic JSON deserialization.
     */
    public PatientCareAdmissionPraxis() {
        super();
    }
}
\end{lstlisting}

% =============================================================================
% SECTION D.4: DYNAMIC ROUTE GENERATION & CAMEL PIPELINE TOPOLOGY
% =============================================================================
\section{Dynamic Route Generation \& Camel Pipeline Topology}
\label{sec:praxis_dynamic_routing}

The core engine of \texttt{PraxisImplementation} is the dynamic route compiler: \texttt{createSequencePipelineRoute()}. When invoked by \texttt{configure()}, it programmatically builds an Apache Camel route tree that seamlessly bridges asynchronous JMS ingress queues, synchronous direct endpoint chaining, pre- and post-activity checkpoint hooks, and completion handoffs.

\subsection{Route Synthesis Architecture}
\label{subsec:routing_synthesis}

Listing~\ref{lst:dynamic_route_synthesis} shows the decompiled dynamic route generation logic employed by \texttt{PraxisImplementation}.

\begin{lstlisting}[language=java, caption={Dynamic Camel Pipeline Synthesis in PraxisImplementation.java}, label={lst:dynamic_route_synthesis}]
@Override
public void configure() throws Exception {
    if (!isEnabled()) {
        LOG.warn("Praxis sequence [{}] is disabled. Skipping route creation.", getPraxisId());
        return;
    }

    // 1. Ingress Conduit Bridge Route: JMS Queue -> Pipeline Start Direct Endpoint
    if (getSourceQueueName() != null && !getSourceQueueName().isBlank()) {
        String jmsUri = getSourceQueueName().startsWith("jms:") 
                ? getSourceQueueName() 
                : "jms:queue:" + getSourceQueueName();
        
        from(jmsUri)
            .routeId("praxis-" + getPraxisId() + "-ingress-bridge")
            .setHeader("PraxisSequenceId", constant(getPraxisId()))
            .to(getPipelineInputEndpoint());
    }

    // 2. Linear Activity Pipeline Chaining
    List<ErgonBase> orderedActivities = getOrderedActivities();
    String startEndpoint = getPipelineInputEndpoint();

    RouteDefinition route = from(startEndpoint)
        .routeId("praxis-" + getPraxisId() + "-pipeline")
        .process(exchange -> {
            Pragma pragma = resolvePragma(exchange);
            if (pragma != null && checkpointManager != null) {
                checkpointManager.recordPipelineIngress(pragma, getPraxisId());
            }
        });

    for (int i = 0; i < orderedActivities.size(); i++) {
        final int stepIndex = i;
        final ErgonBase activity = orderedActivities.get(i);
        String stepEndpoint = "direct:seq-" + getPraxisId() + "-step-" + stepIndex;

        route = route
            .to(stepEndpoint)
            .process(exchange -> {
                Pragma pragma = resolvePragma(exchange);
                if (pragma != null && checkpointManager != null) {
                    checkpointManager.recordPreErgon(pragma, getPraxisId(), stepIndex, activity.getErgonId());
                }
            })
            .to(activity.getInputEndpointUri())
            .process(exchange -> {
                Pragma pragma = resolvePragma(exchange);
                if (pragma != null && checkpointManager != null) {
                    checkpointManager.recordPostErgon(pragma, getPraxisId(), stepIndex, activity.getErgonId());
                }
            });
    }

    // 3. Pipeline Completion Handoff
    route
        .process(exchange -> {
            Pragma pragma = resolvePragma(exchange);
            if (pragma != null && checkpointManager != null) {
                checkpointManager.recordPipelineComplete(pragma, getPraxisId(), orderedActivities.size());
            }
            exchange.getMessage().setHeader("PraxisExecutionStatus", "COMPLETED");
        })
        .to(getPipelineCompleteEndpoint());
}
\end{lstlisting}

\subsection{Exchange Property Invariants and Header Preservation}
\label{subsec:routing_invariants}

Throughout the lifecycle of a Praxis execution pipeline, Apache Camel exchanges must preserve specific headers and properties to maintain transaction traceability. Table~\ref{tab:praxis_exchange_properties} lists the mandatory exchange parameters.

\begin{table}[htbp]
\centering
\small
\begin{tabularx}{\textwidth}{l l p{2.2cm} X}
\toprule
\textbf{Property / Header Name} & \textbf{Scope} & \textbf{Type} & \textbf{Description \& Platform Usage} \\
\midrule
\texttt{PraxisSequenceId} & Header & \texttt{String} & The unique identifier of the executing Praxis sequence. \\
\texttt{ErgonBase.PROPERTY\_PRAGMA} & Property & \texttt{Pragma} & The canonical in-flight task envelope containing clinical payloads. \\
\texttt{Exchange.CORRELATION\_ID} & Header & \texttt{String} & End-to-end correlation ID linking ingress MLLP to outbound dispatch. \\
\texttt{PraxisExecutionStatus} & Header & \texttt{String} & Execution state indicator (\texttt{PROCESSING}, \texttt{COMPLETED}, \texttt{FAILED}). \\
\texttt{PraxisStepIndex} & Property & \texttt{Integer} & The zero-based index of the currently executing pipeline step. \\
\bottomrule
\end{tabularx}
\caption{Mandatory Apache Camel Exchange Properties and Headers in Praxis Pipelines}
\label{tab:praxis_exchange_properties}
\end{table}

% =============================================================================
% SECTION D.5: DYNAMIC DISCOVERY, SEEDING & RUNTIME REGISTRATION
% =============================================================================
\section{Dynamic Discovery, Seeding \& Runtime Registration}
\label{sec:praxis_discovery_seeding}

Harmonia enables zero-downtime reconfiguration of workflow pipelines through \texttt{TaskSequenceLoader} and \texttt{PraxisService}.

\subsection{CDI Bean Resolution via Jakarta EE Instance}
\label{subsec:discovery_cdi}

When loading declarative definitions from the \texttt{tasksequence-cache}, \texttt{TaskSequenceLoader} parses the configured \texttt{activityClassNames} and uses Jakarta EE's \texttt{Instance<ErgonBase>} to dynamically discover and inject matching CDI beans:

\begin{lstlisting}[language=java, caption={Dynamic CDI Ergon Resolution in TaskSequenceLoader.java}]
@ApplicationScoped
public class TaskSequenceLoader {

    @Inject
    private PraxisService praxisService;

    @Inject
    private Instance<ErgonBase> ergonBeanSource;

    @Inject
    private CamelContext camelContext;

    public PraxisImplementation resolveAndInstantiate(PraxisDefinition definition) {
        PraxisImplementation impl = new PraxisImplementation(definition);
        Map<Integer, ErgonBase> resolvedActivities = new HashMap<>();

        for (Map.Entry<Integer, String> entry : definition.getActivityClassNames().entrySet()) {
            int stepIndex = entry.getKey();
            String className = entry.getValue();

            try {
                Class<?> activityClass = Class.forName(className);
                Instance<?> childInstance = ergonBeanSource.select(activityClass);
                
                if (!childInstance.isResolvable()) {
                    throw new IllegalStateException("CDI unable to resolve Ergon bean: " + className);
                }

                ErgonBase ergonBean = (ErgonBase) childInstance.get();
                resolvedActivities.put(stepIndex, ergonBean);
            } catch (ClassNotFoundException e) {
                throw new IllegalArgumentException("Configured Ergon class not found on classpath: " + className, e);
            }
        }

        impl.setActivities(resolvedActivities);
        impl.setCamelContext(camelContext);
        return impl;
    }
}
\end{lstlisting}

\subsection{Initial Seeding via TaskSequenceDefaultSeeder}
\label{subsec:discovery_seeding}

During initial platform boot, the \texttt{TaskSequenceDefaultSeeder} inspects the \texttt{tasksequence-cache}. If the cache is empty, it populates the baseline clinical workflows:
\begin{itemize}
    \item \texttt{praxis-patient-care-admission}: Inpatient and emergency admission pipeline.
    \item \texttt{praxis-lab-order-routing}: Laboratory diagnostic order routing and distribution.
    \item \texttt{praxis-observation-enrichment}: Pathology result enrichment and LOINC terminology mapping.
\end{itemize}

% =============================================================================
% SECTION D.6: MULTI-TIER TESTING STRATEGIES
% =============================================================================
\section{Multi-Tier Testing Strategies}
\label{sec:praxis_testing}

To guarantee pipeline determinism and fault isolation, workflow sequences must be verified using unit, integration, and failure snapshot testing harnesses.

\subsection{CamelTestSupport Unit Test Harness}
\label{subsec:testing_camel}

Listing~\ref{lst:praxis_test_harness} illustrates an automated unit test leveraging Apache Camel Test Kit to verify linear execution, intermediate step transformations, and mock assertions.

\begin{lstlisting}[language=java, caption={Praxis Dynamic Pipeline Unit Test with Mock Assertions}, label={lst:praxis_test_harness}]
package net.fhirfactory.harmonia.praxis.sequence;

import net.fhirfactory.harmonia.erga.base.ErgonBase;
import net.fhirfactory.harmonia.model.pragma.Pragma;
import net.fhirfactory.harmonia.model.pragma.PragmaCheckpoint;
import net.fhirfactory.harmonia.model.pragma.PragmaStatus;
import org.apache.camel.Exchange;
import org.apache.camel.builder.RouteBuilder;
import org.apache.camel.component.mock.MockEndpoint;
import org.apache.camel.test.junit5.CamelTestSupport;
import org.junit.jupiter.api.BeforeEach;
import org.junit.jupiter.api.DisplayName;
import org.junit.jupiter.api.Test;

import java.util.HashMap;
import java.util.List;
import java.util.Map;

import static org.junit.jupiter.api.Assertions.*;

public class PraxisWorkflowUnitTest extends CamelTestSupport {

    private PraxisImplementation praxis;
    private MockErgonStep0 step0;
    private MockErgonStep1 step1;

    @BeforeEach
    public void setupPipeline() throws Exception {
        step0 = new MockErgonStep0();
        step1 = new MockErgonStep1();

        context.addRoutes(step0);
        context.addRoutes(step1);

        praxis = new PraxisImplementation();
        praxis.setPraxisId("praxis-unit-test");
        praxis.setEnabled(true);
        praxis.setCamelContext(context);

        Map<Integer, ErgonBase> activities = new HashMap<>();
        activities.put(0, step0);
        activities.put(1, step1);
        praxis.setActivities(activities);

        context.addRoutes(praxis);
        context.addRoutes(new RouteBuilder() {
            @Override
            public void configure() {
                from(praxis.getPipelineCompleteEndpoint()).to("mock:pipeline-complete");
            }
        });
    }

    @Test
    @DisplayName("Verify complete linear execution of 2-step Praxis pipeline")
    public void testSuccessfulPipelineExecution() throws InterruptedException {
        MockEndpoint completeMock = getMockEndpoint("mock:pipeline-complete");
        completeMock.expectedMessageCount(1);

        Pragma inputPragma = new Pragma();
        inputPragma.setPragmaId("PRAGMA-TEST-001");
        inputPragma.setStatus(PragmaStatus.QUEUED);

        template.sendBody("direct:seq-praxis-unit-test-start", inputPragma);

        completeMock.assertIsSatisfied();
        Exchange resultExchange = completeMock.getReceivedExchanges().get(0);
        Pragma resultPragma = resultExchange.getIn().getBody(Pragma.class);

        assertNotNull(resultPragma);
        assertEquals("COMPLETED", resultExchange.getMessage().getHeader("PraxisExecutionStatus"));
    }
}
\end{lstlisting}

\subsection{Failure Snapshotting and Checkpoint Assertions}
\label{subsec:testing_failure_snapshots}

When an intermediary \texttt{ErgonBase} activity fails (e.g., due to an unrecoverable database outage or payload validation exception), tests must verify that:
\begin{enumerate}
    \item The Camel error boundary intercepts the exception.
    \item A \texttt{PragmaCheckpoint} with stage \texttt{STAGE\_PIPELINE\_FAILED} is recorded in the \texttt{pragma-cache}.
    \item The in-flight \texttt{Pragma} status is transitioned to \texttt{FAILED} with complete error metadata.
\end{enumerate}

% =============================================================================
% SECTION D.7: DEPLOYMENT, CLUSTER MANAGEMENT & OPERATIONAL RUNBOOK
% =============================================================================
\section{Deployment, Cluster Management \& Operational Runbook}
\label{sec:praxis_ops_runbook}

\subsection{Containerization within Ponos WildFly Container}
\label{subsec:ops_containerization}

Praxis workflow sequences run inside the \texttt{hie-task-processor} container (WildFly 31, Java 21). The deployment descriptor (\texttt{pom.xml}) packages compiled Praxis and Ergon classes into the root WAR or shared EAR library.

\subsection{REST Administrative API: /workflow}
\label{subsec:ops_rest_api}

The \texttt{WorkflowManagementResource} exposes administrative endpoints on port 8080:
\begin{itemize}
    \item \textbf{\texttt{GET /workflow/status}}: Returns operational status of all registered Praxis pipelines, including sequence ID, active state, activity count, and target gateway bindings.
    \item \textbf{\texttt{GET /workflow/validate}}: Executes topological validation across all loaded sequences, reporting configuration errors or missing CDI Ergon beans.
    \item \textbf{\texttt{POST /workflow/reload}}: Triggers \texttt{TaskSequenceLoader.reloadAll()}, synchronizing routes with the latest \texttt{tasksequence-cache} definitions.
\end{itemize}

\subsection{Dynamic Live Reload via ponos-cli}
\label{subsec:ops_cli_reload}

Platform operators perform live sequence reloads without restarting the JVM container using the Hestia Operations CLI:

\begin{lstlisting}[language=bash, caption={Operating Dynamic Sequence Reloads with ponos-cli}]
# Check health and status of active Praxis pipelines
ponos-cli workflow status --host ponos-node-1 --port 8080

# Validate declarative configuration in tasksequence-cache
ponos-cli workflow validate

# Execute cluster-wide dynamic pipeline route compilation
ponos-cli workflow reload --force

# Inspect detailed checkpoint trail for a specific task
hestia-cli pragma inspect --pragma-id PRAGMA-992-B71
\end{lstlisting}

% =============================================================================
% SECTION D.8: BEST PRACTICES CHECKLIST
% =============================================================================
\section{Summary \& Best Practices Checklist}
\label{sec:praxis_checklist}

Table~\ref{tab:praxis_readiness_checklist} summarizes the mandatory engineering practices and verification criteria for deploying new Praxis workflow sequences into production.

\begin{table}[htbp]
\centering
\small
\begin{tabularx}{\textwidth}{l p{4.2cm} X}
\toprule
\textbf{Lifecycle Stage} & \textbf{Mandatory Requirement} & \textbf{Production Verification Criteria} \\
\midrule
\textbf{Definition} & Decouple Metadata from Runtime & Use \texttt{PraxisDefinition} for JSON serialization and \texttt{PraxisImplementation} for Camel execution. \\
\textbf{Indexing} & Zero-Based Activity Indices & Ensure \texttt{activityClassNames} and \texttt{activities} maps are indexed sequentially starting at \texttt{0}. \\
\textbf{CDI Injection} & Scoped Ergon Resolution & Annotate activity modules with \texttt{@Dependent} to ensure clean per-route bean isolation. \\
\textbf{Durability} & Multi-Stage Checkpointing & Verify \texttt{PraxisCheckpointManager} writes \texttt{PIPELINE\_INGRESS}, \texttt{PRE\_ERGON}, and \texttt{POST\_ERGON} audit records. \\
\textbf{Resilience} & Dedicated DLQ Routing & Bind pipeline failure boundaries to \texttt{petasos.queue.dlq.workflow-failure}. \\
\textbf{Testing} & CamelTestSupport Mock Suite & Validate happy path, intermediate checkpoint generation, and failure snapshotting. \\
\textbf{Operations} & REST \& CLI Hot Reload & Validate \texttt{/workflow/reload} and \texttt{ponos-cli} synchronization without container restart. \\
\bottomrule
\end{tabularx}
\caption{Production Readiness Checklist for Praxis Workflow Sequences}
\label{tab:praxis_readiness_checklist}
\end{table}
